The residual-finiteness question has several equivalent class-wide formulations. Ilya Kapovich and Wise proved that all hyperbolic groups are residually finite if and only if every nontrivial hyperbolic group has a nontrivial finite quotient \cite[Theorem~1.2]{KW2000}; they also proved equivalence with the assertion that every hyperbolic group is virtually torsion-free, that is, has a torsion-free subgroup of finite index \cite[Theorem~5.1]{KW2000}. Together with Theorem~\ref{thm:main}, these equivalences imply that some nontrivial hyperbolic group has no nontrivial finite quotient, and that some hyperbolic group is not virtually torsion-free. These are conclusions about the class, not additional properties of the torsion-free group constructed here.

Agol, Groves and Manning established a further consequence of a positive answer. A subgroup is \emph{separable} if every element outside it can be distinguished from it in some finite quotient, and a subgroup of a hyperbolic group is \emph{quasiconvex} if geodesics joining its elements remain within a bounded distance of the subgroup. They proved that, if all hyperbolic groups were residually finite, then every quasiconvex subgroup of every hyperbolic group would be separable \cite[Theorem~0.1]{AGM2009}. Their filling theorem produces hyperbolic quotients that simplify intersections among conjugates of the subgroup while keeping a prescribed element outside its image \cite[Theorem~0.6]{AGM2009}. This explains the universal hypothesis: the argument needs residual finiteness of the quotient groups, not just of the initial group.

Additional geometric structure does force residual finiteness. Haglund and Wise introduced special cube complexes and proved that fundamental groups of compact virtually special cube complexes are linear; in particular, these groups are residually finite \cite[Theorem~4.4]{HW2008}. Agol proved that a word-hyperbolic group acting properly and cocompactly on a CAT(0) cube complex is virtually special \cite[Theorem~1.1]{Agol2013}. Thus a hyperbolic counterexample cannot admit such a cubulation. Strong angular-link bounds on a finite polygonal complex do not by themselves provide the cubical structure required by this theorem.

The hyperbolicity hypothesis in Agol's theorem is essential. Wise's complete square complexes, constructed in his 1996 thesis and published in \cite{Wise2007}, include compact nonpositively curved examples whose fundamental groups are not residually finite and whose universal covers are products of trees \cite[Theorem~3.8 and Main Theorem~7.5]{Wise2007}. Burger and Mozes constructed infinite finitely presented simple groups acting freely and cocompactly on products of trees \cite{BM1997}. These tree products contain Euclidean planes and are not hyperbolic. They therefore demonstrate the compatibility of nonpositive curvature with failure of residual finiteness, but not the compatibility with word hyperbolicity established here.
